%=====================================================================
% _     _ _      ______           _
%| |   (_) |     | ___ \         (_)
%| |    _| |_    | |_/ /_____   ___  _____      __
%| |   | | __|   |    // _ \ \ / / |/ _ \ \ /\ / /
%| |___| | |_ _  | |\ \  __/\ V /| |  __/\ V  V /
%\_____/_|\__(_) \_| \_\___| \_/ |_|\___| \_/\_/
%=====================================================================
%interwind
%interwoven
 \chapter{Literature Review on the Modeling of Breakage}
\label{cha:literature_review}
\fancyhead[ER]{Literature Review on the Modeling of Breakage}
\fancyhead[OL]{Literature Review on the Modeling of Breakage}
%=====================================================================
%=====================================================================

The deagglomeration of particles in fluid flows has been the subject
of research for a long time. Various modeling approaches with varying
degrees of complexity are found in the literature.  Provided that the
presence of agglomerates is a prerequisite for breakage, the first
section of this chapter looks at the available methods for handling
the complex structure of agglomerates. Afterwards, studies focusing on
the breakup of agglomerates due to fluid stresses and wall impact are
explored.  The aim is to assess the state of existing models and to
identify areas for potential improvements. Note that major parts of
this review were published in the papers related to this thesis
\citep{breuer2019revisiting,breuer2019refinement,khalifa2020data,khalifa2021efficient,khalifa2021les,khalifa2022neural}.
 \newline

It must be mentioned that breakage might additionally arise due to
collisions between agglomerates and other agglomerates or
particles. However, this mechanism is not addressed in the present
thesis. Therefore, the collisions-induced breakage is not covered in
this literature review.

%That
%is especially the case due to the difficulties encountered in
%obtaining detailed experimental measurements, such as the wide range
%of relevant
%scales and the restricted optical accessibility due to the presence
%of particles.
%=====================================================================
%=====================================================================
\section{Structure of Agglomerates}
\label{sec:litrev_structure}
\fancyhead[OL]{Structure of Agglomerates}
%=====================================================================
%=====================================================================

Agglomerates, also commonly denoted aggregates or flocs, describe
clusters of particles interconnected by different types of binding
mechanisms \citep{fayed2013handbook}. In some cases the geometry of
agglomerates can be highly irregular, whereas in other cases
agglomerates are referred to as fractal-like objects, since they show
a high degree of self-similarity, but are not infinitely
scale-invariant \citep{friedlander1977}. Hence, the morphology of
agglomerates is typically characterized by a parameter denoted the
fractal dimension, which varies between unity for chain-like
arrangements and three for spherical ones. Exemplary
sphere-like and open agglomerates are depicted in
Figs.~\ref{subfig:sphere_like_agg}~and~\ref{subfig:open_agg},
respectively.

A special case is the isostatic structure
(\cref{subfig:open_iso_agg}),  which exhibits a number of
inter-particle contacts less than or equal to
the number of primary particles minus one. Such a morphology evolves
when chains
of particles are cross-linked to form a larger agglomerate.  The
assumption of isostatic  agglomerates is frequently encountered in
the literature, since it  simplifies the modeling
of breakage as will be discussed in
\cref{sec:break_fluid_rigid}. \newline

Particle-laden flow simulation methodologies which deal with cohesive
particles can be distinguished into structure resolving and structure
idealizing methods. In the former, the detailed morphology of
agglomerates is fully taken into account
(e.g.,~Figs.~\ref{subfig:sphere_like_agg}~--~\ref{subfig:open_iso_agg}),
whereas the latter methods replace the structure by a single effective
particle (e.g., Fig.~\ref{subfig:effective_agg}). Both techniques are
discussed in the following.
%
%---------------------------------------------------------------------
\begin{figure}
	\label{fig:particle_laden_exam}
	\centering
	\psfrag{dag}[Bc][B][1][0]{$d_{\text{ag}}$}
    \psfrag{rho}[Bc][B][1][0]{$\rho_{\text{ag}}$}
	%
	\subfloat[Sphere-like structure
	\label{subfig:sphere_like_agg}]{
		\includegraphics[draft=false,width=0.23\linewidth,trim=
		0.0cm 0cm 0cm
		0.0cm,clip]{figure/lit_rev/sphere_like_agg.eps}}%
	\hfill
	%
	\subfloat[Open structure
	\label{subfig:open_agg}]{
		\includegraphics[draft=false,width=0.23\linewidth,trim=
		0.0cm 0.0cm 0.0cm 0.0cm,clip]{figure/lit_rev/open_agg.eps}}%
	\hfill
	%
	\subfloat[Open, isostatic structure
	\label{subfig:open_iso_agg}]{
		\includegraphics[draft=false,width=0.23\linewidth,trim=
		0.0cm 0.0cm 0.0cm
		0.0cm,clip]{figure/lit_rev/open_iso_agg.eps}}%
		\hfill
	%
	\subfloat[Effective sphere
	\label{subfig:effective_agg}]{
		\includegraphics[draft=false,width=0.23\linewidth,trim=
		0.0cm 0.0cm 0.0cm
		0.0cm,clip]{figure/lit_rev/effective_agg.eps}}%
	\\ \vspace{0.5cm}
	\caption{Agglomerates represented by  detailed
			structures (a)-(c) and
			an idealized one based on an effective sphere (d).}
\end{figure}
%\texttt{http://visibleearth.nasa.gov}
%---------------------------------------------------------------------

%=====================================================================
\subsection{Detailed Structure}
%=====================================================================

 In the first category the structure of the agglomerate is fully taken
 into account, i.e., agglomerates are realized as assemblies of
 primary particles. In addition, the dynamics of the agglomerates are
 determined by the fluid-particle and inter-particle interactions
 occurring at the scale of the constituting primary particles. Hence,
 both types of interactions need to be adequately addressed.



 The \emph{free-draining} approximation is the simplest approach
 suggesting to entirely ignore the fluid shielding effects
 \citep[e.g.,][]{becker2009restructuring}.  In a more reasonable
 technique, the effect of the neighboring particles is accounted for
 by incorporating a \emph{crowding factor} into the point-particle
 force and torque models. For this purpose, a large number of
 experimental \citep[e.g.,][]{di1994voidage} and particle-resolved
 numerical \citep[e.g.,][]{beetstra2007drag} studies proposed
 correlations describing the crowding factor as a function of the
 particle Reynolds number and the mean local volume fraction of the
 particles. However, most correlations are mainly focused on the drag
 force and independent of the effect of the spatial arrangement (i.e.,
 relative locations) of the particles, which can strongly vary despite
 an equal local volume fraction.

 To overcome these drawbacks, different possibilities are
 available. One of them is the \emph{hydrodynamic screening} method,
 which computes the fluid-induced forces and torques exerted on a
 particle by considering solely the portion of the surface area of the
 particle exposed to the flow field \citep[]{higashitani2001}.
 Another alternative is the \emph{Stokesian dynamics}
 \citep{durlofsky1987dynamic}, which allows to calculate the fluid
 forces and torques exerted on the primary particles taking the
 relative positions of particles into account. In this method the
 fluid forces and moments are related to the velocities of the
 particles by the so-called mobility matrix, assuming that particles
 are surrounded by a flow field with a uniform velocity gradient,
 i.e., a creeping flow. A more generally applicable technique is the
 \emph{pairwise-interaction extended point-particle (PIEPP)} model,
 which was recently proposed by
 \citet{akiki2017pairwise,akiki2017pairwisepar,balachandar2020towards}.
 The model allows to compute the drag and lateral forces on individual
 particles within a random array while systematically accounting for
 realistic surrounding flow conditions and for the relative positions
 of the neighboring particles. \newline

In principle, all above mentioned approaches can be considered for
describing the fluid forces acting on aggregated particles while
dealing with the inter-particle interactions based on one of two
typical options: (1) the discrete element method (DEM), (2) the
rigid-body dynamics.

In the present terminology, \emph{DEM} refers to the method, in which
the inter-particle forces and torques are computed based on a
soft-sphere (i.e., viscoelastic) model. In addition, the velocity and
displacement of each individual particle are obtained by integrating
its equation of motion taking the external (i.e., fluid and gravity)
and the internal (i.e., contact) forces into account
\citep{cundall1979}. Hence, in DEM the attachment and detachment of
particles are directly dictated by the resolved motion of
particles. However, a stable and accurate DEM solution requires very
small time increments, which can be much smaller than the flow
time-step size especially for stiff particles (i.e., high modulus of
elasticity).  Therefore, it is a very common practice to artificially
scale down the modulus of elasticity to allow larger DEM time-step
sizes. Note that DEM studies adopting the direct prediction for the
breakup by fluid stresses and by wall impacts are discussed in
Sections~\ref{sec:break_fluid_DEM}~and~\ref{sec:break_wall_DEM},
respectively.

The second possibility is more efficient since it treats agglomerates
as \emph{rigid structures}
\citep[e.g.,][]{seto2011hydrodynamic,vanni2011hydrodynamic,fellay2012effect,schutte2015hydrodynamic,debona2014}.
In practice, rigid agglomerates indicate stable inter-particle bonds
which remain fixed or spontaneously break when the exerted forces or
torques exceed certain thresholds. Computationally, the motion of a
rigid agglomerate is described by the translational and rotational
velocities of the center of mass of the agglomerate. These velocities
are obtained by solving the corresponding equations of motion of the
whole agglomerate taking the external forces and torques acting on the
constituting primary particles into account. The main advantage of
this approach over typical DEM is that it does not necessitate the
prediction of the internal interactions or solving the equations of
motion for each primary particle.  However, the assumed rigidity of
agglomerates implies that the primary particles are always intact and
thus additional criteria are required to determine breakage. Further
discussion on the modeling of the fluid-induced breakage for rigid
agglomerates is given in Section~\ref{sec:break_fluid_rigid}. Note
that to the best of the author's knowledge, studies dealing with the
wall-impact breakage of rigid agglomerates with detailed structures do
not exist in the literature.
%by a two-step
%procedure. In the first step, the total perturbation flow caused by
%the presence of neighboring particles is obtained by superimposing
%the perturbation flow due to each particle separately.

%=====================================================================
\subsection{Idealized Structure \label{sec:idealizied_stru}}
%=====================================================================

In the second category of studies, the detailed agglomerate morphology
is replaced by a single particle conserving the mass of the original
agglomerate. Very often the model particle is assumed to be a sphere
characterized by an effective diameter (see
Fig.~\ref{subfig:effective_agg}).

Due to its practicality, the
\emph{replacement technique} has been extensively utilized in the
literature. For instance, the effective agglomerate diameter is an
essential parameter which frequently appears in the closures, e.g.,
the agglomeration and breakage kernels used in the population balance
equation. The replacement approach is also popular in the
Euler--Lagrange formalism
\citep[e.g.,][]{breuer2015a,kosinski2011,ho2002,chun2005coagulation,van2020analysis}
since it offers a considerable reduction of the computational load by
giving up the continuous tracking and updating of the inter-particle
contacts. In addition, modeling agglomerates as single spherical
entities allows to apply the same closure relationships for the fluid
forces (e.g., drag and lift and their corresponding drag and lift
coefficients) developed for spherical primary particles to predict the
motion of agglomerates. The effective sphere model is especially
beneficial for simulations parallelized by domain decomposition
techniques since the inter-processor communication and exchange of
aggregated particles is not a simple task. \newline

Certainly, an effective sphere provides solely an approximation of the
dynamics of the original agglomerate.  The achieved accuracy was found
to be affected by the considered definition of the effective diameter
\citep{dietzel2016application}. In most studies, a volume-equivalent
diameter is taken into account, which conserves the total volume as
well as the total mass of the primary particles.  That requires the
density of the agglomerate to be equal to that of the primary
particles ignoring the interstitial spaces between the particles.

\citet{breuer2015a} followed an extended approach denoted
closely-packed sphere model. It accounts for the porous nature of
agglomerates by defining the density of an agglomerate as the product
of the density of the primary particle and a prescribed porosity
value.  Consequently, the volume of the sphere representing the
agglomerate structure is larger than the total volume of the
constituting particles. Since agglomerates of different number of
primary particles might arise in the simulations, the packing fraction
(i.e., the complement of the porosity) of agglomerates was provided in
a look-table depending on the number of particles.  The tabulated
values of the packing fraction stem from an optimization study on the
smallest volume of a cube enclosing a certain number of equal spheres
without overlapping \citep{packomania2013}. Due to the assumption of
closely packed particles in a cube, the resulting values of the
packing fraction are found to be very high in comparison to the
typical range of values for randomly packed, spherical particles
\citep[e.g.,][]{yang2008agglomeration,scott1962radial,finney1970random}.\newline
% \newline
%Another drawback of the approach is that it only accounts for the
%number of primary particles constituting the agglomerate in
%determining its density and thus effective size.  Other relevant
%factors, such as the packing method and the material properties,
%namely cohesion are neglected.
%
%---------------------------------------------------------------------
\begin{figure}
	\centering
	\psfrag{hull}[BL][Bc][1.][0]{Convex hull}
	\includegraphics[draft=false,width=0.4\textwidth]{figure/lit_rev/convex_tecplot.eps}%
	\caption{Convex hull wrapping an agglomerate consisting of 100
		primary particles.}
	\label{fig:convex_hull}
\end{figure}
%---------------------------------------------------------------------
%

To assess the suitability of the effective sphere assumption,
\citet{harshe2010hydrodynamic} simulated the motion of rigid
agglomerates in a uniform shear flow by means of Stokesian
dynamics. Agglomerates of sizes varying between 10 and 1000 primary
particles, possessing different open and sphere-like structures were
investigated. It was shown that the simplifying effective sphere
approach is useful to reproduce the behavior of realistic
agglomerates. Specifically, sphere-like agglomerates can be well
represented by equivalent spheres, whereas ellipsoids are more suited
for open structures.

Another important contribution in this area was presented by
\citet{dietzel2013numerical} who performed particle-resolved
lattice-Boltzmann simulations at low Reynolds numbers (i.e., Stokes
regime) to investigate drag and lift force coefficients of
agglomerates possessing different open and sphere-like structures. It
was concluded that a drag force coefficient computed based on the
diameter of a volume-equivalent sphere is not adequate since it
remarkably deviates from the standard drag coefficient correlation by
\citet{schiller33}.  A considerably better agreement was achieved
based on a sphere with an area equal to the projected surface area of
the convex hull wrapping the agglomerate (see
Fig.~\ref{fig:convex_hull}). Namely, the diameter of a circle
possessing an area equal to the projected cross-section area of the
convex hull resulted in drag coefficients which are nearly independent
of the morphology and the orientation of the agglomerates
\citep{dietzel2013numerical,dietzel2016application}. However, unlike
the drag coefficient, the lift force coefficient was not well
predicted by any of the two structural assumptions (sphere and convex
hull).

In a subsequent study, \citet{sommerfeld2017novel} proposed a
strategy for predicting the effective diameters of agglomerates based
on a volume-equivalent convex hull. The model was developed for
Euler-Lagrange point-particle simulations and can be summarized as
follows: When an agglomerate appears, the volume of the convex hull
enclosing the aggregated particles is computed. Subsequently, the
relative positions of the agglomerated particles are stored and the
agglomerate is replaced by a sphere with a volume equal to that of the
convex hull. Upon further agglomeration or breakup, the positions of
the added or removed particles are used together with the stored
relative positions of the current particles to compute the volume of a
new convex hull, and so forth. \newline

The breakup of agglomerates idealized as spheres due to the
interactions with a fluid flow field and as a result of wall impacts
are discussed next in
Sections~\ref{sec:break_fluid_sphere}~and~\ref{sec:break_wall_replace},
respectively.

%=====================================================================
%=====================================================================
\section{Breakup by Fluid Stresses}
\label{sec:litrev_fluidstress}
\fancyhead[OL]{Breakup by the fluid stresses}
%=====================================================================
%=====================================================================

The role of fluid forces in splitting agglomerates has been
computationally studied in a growing body of literature. Loosely
speaking, independent from the level of resolution of the
fluid-particle coupling (particle-resolved or point-particle) the
breakup of agglomerates can be reproduced based on
\emph{condition-free} methods or by imposing {external breakup
  conditions}. Note that "condition-free" and "imposed conditions" in
this context strictly refer to the forecasting of breakup.

%=====================================================================
\subsection{Fluid-Induced Breakup Predicted by Condition-Free
Techniques}
\label{sec:break_fluid_DEM}
%=====================================================================

As aforementioned, to enable a condition-free forecast of
fragmentation of agglomerates the detailed structure must be taken
into account. In addition, the equations of motion of each primary
particle need to be solved.  Consequently, breakage follows directly
from the equations of motion of the individual particles without
additional modeling. Hence, in all studies discussed in this section
the structures of agglomerates are fully taken into account and the
inter-particle interactions are treated according to DEM.

 Studies utilizing the most accurate approach, which combines
 \textbf{resolved
 fluid-particle interactions} with DEM have been concerned with the
 development of dry powder inhalers
 \citep[e.g.,][]{cui2018application,vanwachem2017simulation}. Within
 the context of inhalers, different powder formulations are usually
 discussed: The carrier-based formulation referring to drug particles
 attached to a larger carrier as depicted in
 Fig.~\ref{fig:carrier_agg}, and the drug-based (carrier-free)
 formulation representing agglomerates of small drug particles
 analogous to those shown in
 Figs.~\ref{subfig:sphere_like_agg}~to~\ref{subfig:open_iso_agg}
 \citep{longest2013aerodynamic}. Bearing in mind that the present
 thesis deals with agglomerates resembling the carrier-free powder
 type, the remainder of this review will focus on studies describing
 the breakage of such agglomerates. Publications addressing the
 disintegration of carrier-based agglomerates are of little value for
 the scope of the present thesis and thus the interested reader is
 referred to the relevant reviews, e.g., by
 \citet{zheng2021flow,yang2015numerical,sommerfeld2019potential}.
 \newline
%Note
%that such an approach is computationally not feasible and is limited
%to fundamental investigations taking mostly a single agglomerate
%into
%consideration.

 %---------------------------------------------------------------------
\begin{figure}
	\centering
	\includegraphics[draft=false,width=0.3\textwidth]{figure/lit_rev/carrier.eps}%
	\caption{Carrier-based powder formulation. This type of
	agglomerates is not considered in the present thesis.}
	\label{fig:carrier_agg}
\end{figure}
%---------------------------------------------------------------------
%
%----Point particle: DEM with crowding factor

Other studies considered \textbf{unresolved particle} methods. These
studies
are grouped in the following based on the approach used for dealing
with the fluid shielding within the agglomerate.

The first group adopted the \emph{hydrodynamic screening method}.  In
one of the earliest contributions, \citet{higashitani2001} studied the
breakup of agglomerates in simple shear and elongational flows.
Agglomerates possessing open, loose structures as well as spherical,
rather compact structures were investigated.  The size of the
agglomerates varied between 256 and 1024 identical particles. The
effect of shielding on the fluid forces experienced by a particle was
accounted for based on a geometrical analysis taking solely the
exposed part of the surface area of the particle into account. In
addition, the undisturbed local flow velocity was adjusted by a factor
accounting for the presence of neighboring particles. It was shown
that elongational flows induce higher breakup than shear flows. In
addition, power-law relations describing the average number of
particles in the broken fragments were proposed. Adopting the same
methodology, \citet{fanelli2006prediction,fanelli2006prediction2}
investigated the breakage of single agglomerates in steady and
unsteady shear flows, respectively. \newline

The second group employed the \emph{Stokesian dynamics}.
\citet{harada2006dependence} studied the fragmentation of agglomerates
in a simple shear flow. The scope was to compare the breakage behavior
of two 100-particle agglomerates distinguished solely by the fractal
dimension (open vs. sphere-like). It was pointed out that
open-structured agglomerates are more fragile than spherical ones and
that the breakup behavior in general depends on the strain and
rotation effects in shear flows.  In an effort to reduce the
computational costs associated with DEM,
\citet{frungieri2021aggregation} presented a method, which combines a
Stokesian dynamics and a stochastic Monte-Carlo algorithm to track
single particles and detect the occurrence of binary collisions. A DEM
approach was incorporated in the framework of Stokesian dynamics to
solely describe the collisions of primary particles and the dynamics
of the formed agglomerates.  The method was applied to study the
agglomeration and breakup under a uniform shear flow condition for a
suspension of monodisperse particles possessing a volume fraction of
10$^{-4}$.  \newline

In the third group of studies, the \emph{free-draining approximation}
was applied.  \citet{becker2009restructuring} investigated the
restructuring of single monodisperse agglomerates in a simple shear
flow applying the free-draining approximation. To assess the
appropriateness of the free-draining assumption, additional finite
element simulations were carried out, in which the fluid flow around a
small agglomerate comprising eight particles was fully resolved.  The
obtained fluid forces on each particle were compared with results of
standard constitutive force models assuming an undisturbed flow field
around the particle. The analysis revealed that the free-draining
assumption leads to an overestimation of the fluid forces experienced
by the particles exposed to the flow field. Furthermore, the
perpendicular forces induced by neighboring particles in the inner
part of the agglomerates are underestimated.

\citet{eggersdorfer2010fragmentation} examined the restructuring of
single agglomerates of identical particles in a simple shear flow.
Agglomerate comprising up to 2048 particles were generated by a
hierarchical stochastic algorithm leading to widely open agglomerate
structures.  It was highlighted that after agglomerates split, the
arising fragments undergo a relaxation phase leading to more compact
structures than the original agglomerates.

  Another study using the
free-draining assumption was published by
\citet{calvert2011mechanistic} who evaluated the effect of cohesion,
the diameter of the agglomerate, and the fluid-agglomerate relative
velocity (commonly referred to as the slip velocity) on the
disintegration of agglomerates in uniform flow. For this purpose,
sphere-like agglomerates comprising a number of primary particles
varying between 250 and 8000 particles were taken into account. In
addition, five different slip velocities were considered and a range
of one order of magnitude of surface energy values was considered. To
characterize the breakup behavior, a parameter called the dispersion
ratio was introduced, which related the number of broken contacts to
the initial number of contacts. It was found that a higher relative
velocity is needed to break an agglomerate if the surface energy of
the particle increases.  In contrary, the onset of breakage for larger
agglomerates takes place at lower slip velocities than in the case of
small agglomerates. A correlation incorporating a dimensionless number
analogous to the Weber number was found to characterize the dependence
of the dispersion ratio and the operating (flow and agglomerate)
conditions.

In a recent study, \citet{zhao2021flocculation}
investigated the dynamics of $10^4$ single cohesive particles in
homogeneous isotropic turbulence employing one-way coupled DNS. A
rapid aggregation phase was initially observed before a steady
equilibrium phase prevailed due to the balance between agglomeration
and breakup. The flow-shielding effects on the agglomerated particles
were neglected.  It was found that agglomerates are elongated along
the strongest Lagrangian stretching direction before breakage. Based
on the obtained results, a
model describing the transient development of the aggregate size and
shape was proposed. \newline

DEM-based methodologies accounting for the flow-shielding effect in
the calculation of the drag force by a \emph{crowding factor} have
been reported in different recent studies (the fourth group).
\citet{dizaji2019collision} investigated the breakup of agglomerates
in a linear shear flow focusing on the effect of the fluid forces and
the collisions of two agglomerates. For this purpose, two sets of
simulations involving one or two agglomerates were carried out.  The
considered agglomerates were generated in a preceding DNS of
homogeneous turbulence and had open (loose) structures consisting of
328 to 577 evenly-sized particles. It was demonstrated that the
breakage behavior, especially for open agglomerate structures, is more
complex than can be reproduced by effective spheres.  Furthermore,
\citet{chen2020collision} applied a coupled DNS--DEM to investigate
the collision-induced breakage of agglomerates in homogeneous
isotropic turbulence. After randomly seeding $4\cdot10^4$ single
primary particles of the same size, the breakage of the small (mostly
two-particle) agglomerates formed during the simulations was
analyzed. In addition, a model for the rate of the collision-induced
breakage in such flow scenarios was proposed. Moreover,
\citet{yao2021deagglomeration} performed two-way coupled DNS--DEM to
investigate the disintegration of single agglomerates in homogeneous
isotropic turbulence. The closely-packed sphere-like agglomerates
comprising 682 and 5500 equally-sized particles were composed by a
centripetal packing procedure prior to the breakup investigations. As
an outcome, a model describing the condition of breakup and the
breakup rate was derived based on arguments motivated by the Taylor
analogy breakup (TAB) model \citep{o1987tab}, which is readily used
for droplet breakup. \newline

Obviously, DEM-based methods are insightful and provide
detailed
knowledge on the primary particle scale. However, to keep the required
computational efforts tractable, simplifications need to be made. Even
in the most recent studies, the scope was restricted to single or few
agglomerates under simple flow conditions. For these reasons, such
detailed methods are still limited to fundamental investigations and
not feasible for practical applications.

%=====================================================================
\subsection{Fluid-Induced  Breakup Predicted by Condition-Based
Techniques}
%=====================================================================

In the conditions-based approaches, the equations of motion of the
primary
particles constituting an agglomerate are not individually solved.
Instead, the agglomerates are tracked either as rigid bodies
consisting of a number of particles, or as effective spheres. In
either case, additional criteria are required to decide upon the
occurrence of breakup. In the case of an effective sphere, additional
models are needed to describe the arising fragment size distribution
and the velocity of the fragments.

%---------------------------------------------------------------------
\subsubsection{Breakup of Rigid Agglomerates with Detailed Structures}
\label{sec:break_fluid_rigid}
%---------------------------------------------------------------------

The first type of the conditions-based methods suggests to track
agglomerates as single rigid entities while accounting for their
detailed structures. This allows to determine the motion of the whole
agglomerate without resolving the inter-particle interactions.
However, in order to formulate a breakup criterion that makes use of
the available morphology, knowledge on the inter-particle forces and
torques becomes necessary. For general agglomerates possessing
arbitrary structures, applying a force or a torque balance (external
forces versus contact forces and likewise for torque) per particle
leads to an ill-defined problem \citep{derksen2010potential}. This
situation arises when the number of unique inter-particle contacts
(i.e., unknowns) is larger than the number of particles (i.e.,
equations).  For this purpose, isostatic agglomerates (e.g.,
Fig.~\ref{subfig:open_iso_agg}) are frequently considered in the
studies dealing with the restructuring (i.e., agglomeration and
breakup) of rigid agglomerates.  An isostatic structure means that the
number of inter-particle contacts is less than or equal to the number
of particles minus one \citep{schutte2018formation}.  As a
consequence, the structure is statically determined, i.e., imposing a
force and a torque balance on each particle yields two closed systems,
which can be solved to obtain the inter-particle forces and torques. 
\newline 

\citet{derksen2010potential} performed \textbf{particle-resolved}
lattice-Boltzmann simulations to study the breakage of agglomerates in
micro-channels under laminar flow conditions.  Agglomerates comprising
solely 2 to 4 particles were tracked as rigid bodies, i.e., the
interactions between the particles of the same agglomerate are not
resolved. However, the collisions between agglomerates were considered
by resolving the interactions between the colliding particles of
different agglomerates based on a soft-sphere interaction model. Since
the agglomerated particles were assumed to maintain their integrity,
the scope of the problem was restricted to analyzing the development
of the normal contact force between the particles. \newline 

Considering an \textbf{unresolved particle} approach on the basis of
\emph{Stokesian dynamics}, \citet{vanni2011hydrodynamic} examined the
behavior of single agglomerates of different fractal dimensions
immersed in simple shear flow. The inter-particle interactions within
the agglomerates were computed by applying a force and a torque
balance on each particle. Since the forecast of breakup requires to
define critical thresholds for the forces and torques in three
directions (i.e., six degrees of freedom), the focus was limited to
examining the evolution of the inter-particle interactions within the
agglomerate. 

 Using the same method, \citet{seto2011hydrodynamic}
studied the relation between the size of the agglomerates and the
forces acting on the particle contacts in simple shear
flows. 

Furthermore, \citet{debona2014} followed a similar approach to
investigate the breakup of agglomerates in homogeneous isotropic
turbulence. The failure of contacts was solely attributed to the
normal (tensile) stresses. Other possible breakup modes such as
shearing, bending and twisting of contacts were not considered.
Consequently, a critical tensile stress model suggested by the contact
mechanics was employed to set a criterion for breakup. A simple case
of two-particle agglomerates and a more complicated one involving
agglomerates possessing open structures were considered.  As an
outcome, a model describing the breakup rate in homogeneous isotropic
turbulence was suggested.  In addition, it was found that the
orientation of the agglomerate with respect to the flow, especially
for the two-particle agglomerate, has considerable effects on breakup.

In a recent work, \citet{frungieri2020cfd} studied the breakup of
rigid aggregates in an internal mixer. Two-dimensional, unsteady
simulations of the flow in the cross-section of the mixer were
conducted at two different rotor speed ratios taking 100 tracer
particles representing agglomerates in a one-way coupled simulation
into account. In parallel, the interactions between the flow field and
the detailed isostatic structure of each agglomerate were computed
applying the Stokesian dynamics method. In addition, the
inter-particle forces and torques were obtained based on balances
imposed at each contact as typical in the rigid-body method. A bond
was assumed to instantaneously rupture when it experiences a tensile
force exceeding a prescribed critical threshold. However, no further
elaboration on the relation between the critical tensile force and the
particle properties was given. 

\citet{schutte2015hydrodynamic} used the rigid, isostatic structure
assumption to study the dynamics of asphaltene agglomerates including
breakup in channel flows. The flow was simulated one-way coupled by
means of DNS and LES considering four turbulent flow conditions and
$2.5 \cdot 10^{5}$ primary particles. Solely the drag and added-mass
forces were considered for the flow-particle interactions and the
flow-shielding effect was neglected (\emph{free-draining
approximation}). Four modes of breakage, i.e., straining (tensile),
shearing, bending, and twisting were considered. However, the values
of the critical forces and torques leading to breakup were varied to
explore the breakup behavior for different levels of bond
strength. Hence, no rigorous derivations of these quantities were
proposed. In a subsequent investigation, \citet{schutte2018formation}
extended their methodology to take the effect of two-way coupling into
account.  The method was applied to study the breakage of agglomerates
in channel flows. It was concluded that agglomerates tend to break
into fragments of asymmetric sizes rather than identical
ones.

Moreover, a simplified version of the rigid-body model exists
\citep{kousaka1979dispersion,kousaka1980re,kousaka1992dispersion,
  niedballa2000modeling}. Therein, agglomerates are idealized by
two-particle (doublet) structures and the contact forces are estimated
to evaluate the possibility of breakup. The main drawback of this
model is that the condition of breakage strongly depends on the
relative orientation of the two-particle structure to the flow
direction and on the size ratio of the two particles. Hence, the
prediction of a breakup is strongly affected by modeling issues.
\newline

Overall, the assumption of rigid, isostatic agglomerates facilitates
less CPU-intensive computations since the motion of each constituting
particle does not need to be directly resolved. However, as pointed
out before, tracking the structure especially in multi-block grids on
parallel computers is rather complicated. Furthermore, isostatic
agglomerates split as soon as a single contact within the agglomerate
breaks up.  This is a specific case which does not necessarily
represent the breakage behavior in the more general case of arbitrary
(hyperstatic) agglomerates.  Finally, the main drawback of this
modeling approach is the lack of rigorous models to describe the
critical forces and torques leading to breakage.

%---------------------------------------------------------------------
\subsubsection{Breakup of Agglomerates Idealized as Spheres}
\label{sec:break_fluid_sphere}
%---------------------------------------------------------------------

The second type of conditions-based methods describes the structure of
agglomerates based on the \emph{effective sphere} approximation.
Numerous studies have adopted this approach to study the fluid-induced
breakup of agglomerates in the context of the Euler-Euler as well as
the Euler--Lagrange framework. In the former, breakage rate models are
required, whereas in the latter beside a breakage condition the size
and the motion of the resulting fragments must be addressed. Such rate
functions and breakage models are typically formulated based on
simplified analytical derivations, experimental observations, or
results of detailed numerical simulations.

%$\mathcal{O} (10^4)$
Characterizing the breakage of agglomerates and deriving predictive
models has been the goal of several theoretical studies
\citep[e.g.,][]{rumpf1932desagglomeration,thomas1964turbulent,
  tomi1978behaviour,sonntag1987structure2,zaccone2009}. Therein, a
general hypothesis was accepted, which states that breakage occurs
when the fluid dynamic stress exceeds a critical threshold that
strongly depends on the type of the inter-particle cohesive forces,
i.e., the yield stress or the strength.  Consequently, the focus
therein has been put on elucidating both sides of the problem
considering agglomerates of idealized structures, e.g., porous spheres
or cylinders, surrounded by various flow patterns ranging from simple
uniform or shear flows to more complicated and irregular turbulent
flows. \newline

Considering \emph{simple flow cases},
\citet{rumpf1932desagglomeration} described different disruptive
mechanisms such as (1) the stresses acting on agglomerates
accelerating or decelerating in a uniform flow field and the (2) shear
and (3) centrifugal stresses imposed on agglomerates rotating in
simple shear flows. The available analytical solutions for the
considered simple flow cases allowed to determine the key flow
quantities responsible for each mechanism.  Based on that, expressions
for computing these fluid-induced stresses were derived considering
various assumptions.  For example, the stress experienced by an
agglomerate accelerating or decelerating in a flow field (often
denoted \emph{inertia stress or drag stress}) was attributed to the
relative velocity between the agglomerate and the fluid flow. Hence,
it has been approximated by the ratio of the drag force acting on a
spherical agglomerate to its surface area. In addition, an expression
for the shear stress acting on the surface of spherical or cylindrical
agglomerates in a shear flow was derived on the basis of the
elasticity theory, assuming agglomerates as elastic objects.
Moreover, the angular velocity of agglomerates in shear flows has been
approximated based on the shear rate of the fluid, so that the
application of the elasticity theory has led to the prediction of the
stress developing in spherical or cylindrical agglomerates due to
rotation, denoted \emph{rotary stress}. Furthermore, the stress
threshold needed for a successful breakup has been introduced based on
the famous model by \citet{rumpf1962}, who has postulated that the
size of the primary particles, their structural properties, and the
type of the binding mechanism are the main factors determining the
strength of agglomerates.

In the seminal work by \citet{tomi1978behaviour} the stresses
developing within a sphere in simple shear or uniform flow fields have
been elaborated in detail aiming to define planes of maximum stress,
which resemble the likely surfaces of failure. The critical stress
needed for breakup was related to the properties of the agglomerates
as suggested by the yield stress model by \citet{thomas1963transport}
which is similar to the model by \citet{rumpf1962}, albeit a slightly
different relationship is used. As an alternative to the model by
\citet{thomas1963transport}, the Bingham plastic yield model was
suggested to estimate the critical yield stress beyond which
agglomerates are assumed to become flowable and therefore
split. \newline

The breakup of agglomerates under \emph{turbulent flow conditions} has
also been studied theoretically and experimentally by several authors.
\citet{thomas1964turbulent} has suggested a breakup model, in which
the mechanism of disruption is mainly referred to the pressure
difference on opposite sides of the agglomerates caused by random
velocity fluctuations. Assuming locally isotropic turbulence,
Kolmogorov's theory of local isotropy has been applied to estimate the
root-mean-square velocity fluctuations over distances comparable to
the diameters of the agglomerates by integrating the energy spectrum
over the range of wave numbers (eddy sizes) related to the size of the
agglomerate. Two main subranges have been considered: The viscous
subrange characterized by eddy sizes smaller than the Kolmogorov's
length scale, and the inertial subrange, in which the eddies are much
larger than the Kolmogorov's length scale but smaller than the
integral length scale of turbulence.  The root-mean-square velocity in
each subrange was used to calculate the turbulent stress imposed on
agglomerates with sizes comparable to the corresponding eddy
size. Consequently, breakage rate functions were derived taking the
critical yield stress model \citep{thomas1963transport} into
account. The rate functions were used to determine an average
agglomerate size denoted steady-state diameter.

Following similar arguments, \citet{neesse1987physikalische} collected
formulas for estimating the root-mean-square velocity fluctuations in
the viscous, inertial, and two intermediate subranges of turbulence
from the literature. Subsequently, the proposed expressions are
applied to compute the turbulent shear or normal stresses based on the
nature of the flow in the corresponding subrange. It has been argued
that the interactions between agglomerates and viscous eddies lead to
surface shearing, and hence the expected breakup mode is erosion,
whereas the interaction with inertial eddies results in bulk
deformations or rupture. Note that these considerations are in line
with the postulations and suggestions of \citet{gregory1989},
\citet{muhle1993coagulation} and \citet{yeung1997effect}.
Accordingly, the strength of the agglomerates has been included in the
model taking the breakup mode (i.e., erosion or bulk rupture) into
account. For rupture, the strength model by \citet{rumpf1962} was
found adequate. However, a pseudo-surface tension model was proposed
to describe strength against erosion \citep{neesse1987physikalische}.
The background of this approach combines ideas from Rumpf's tensile
strength model and from the Bingham plastic yield model.

Motivated by observations of agglomerates breaking down due to
unstable propagation of internal cracks, \citet{zaccone2009} provided
a model, which aims at predicting the largest stable diameter of
agglomerates in laminar and turbulent regimes. Applying an energy
balance, the occurrence of breakup was attributed to the circumstance
that the supplied energy by the fluid dynamic stresses exceeds the
energy required for extending the fracture surface. Employing several
empiric scaling laws and assumptions, a relationship correlating the
applied stress $\sigma$ and the steady-state diameter of the
agglomerate $d_{ag}$ was suggested in the form $\sigma = m \:
d_{ag}^n$. This power-law relationship is generally in agreement with
several experimental findings
\citep[e.g.,][]{adler1979,sonntag1986structure,serra1997aggregation,
  wengeler2007turbulent,soos2008effect,ammar2012,saha2015breakup}.
However, it is noted that the correlation parameters depend on the
flow configuration, the material, and the structure of the
agglomerates, which explains the deviations in the values provided for
the parameters $m$ and $n$ in the aforementioned studies. \newline

The models discussed so far have been used in several subsequent
\emph{computational studies}. For example, \citet{disskusters1991}
studied the agglomeration and breakage of agglomerates comprising
small polystyrene particles (1 $\mu$m) in isotropic and homogeneous
turbulence. The breakup condition was derived by equating the
turbulent stresses as introduced to the strength of the agglomerate
based on an extended version of the model by \citet{rumpf1962} which
accounts for the fractal dimension of the agglomerate. Based on these
considerations a breakup rate function was derived.

\citet{ammar2012} performed an Euler--Euler study on the breakup of
Titanium dioxide agglomerates (0.1--10 $\mu$m) in a highly turbulent
pipe flow (Re = $10^5$--$10^7$). The breakup rate was derived taking 
the
turbulent stresses and the strength of the agglomerates proposed by
Rumpf into account. Despite neglecting its influence, the authors have
pointed out the relative importance of the drag stress mechanism in
their case, since the formed agglomerates are expected to possess
considerable inertial characteristics. \citet{babler2008} investigated
the breakage of small and light agglomerates in homogeneous turbulence
using the Euler--Euler approach. Breakup was assumed to be invoked by
sufficiently strong turbulent stresses. The rate of breakup was
determined using a power-law function describing the maximum stable
size of agglomerates under a given stress. \newline 

%
%Within the Euler--Lagrange framework, the assessment of breakup
%in this technique follows from the evaluation of the fluid stresses
%computed at the position of the point-particle agglomerates along
%their trajectories and a subsequent comparison with a critical
%threshold describing the cohesive strength.

Within the Euler--Lagrange framework, a frequently reported practice
is to record the stresses acting on the point agglomerate to evaluate
the possibility of breakup, without actually modeling the
fragmentation event. According to a very recent review by
\citet{zheng2021flow} on the modeling of the flow in dry powder
inhalers, this approach has been widely applied in the context of
inhalers and is generally referred to as \emph{particle trajectory
  modeling}.  For example,
\citet{weiler2010new,dissweiler2008generierung} computed the
turbulent, drag and rotary stresses acting on agglomerates while they
are transported through an accelerated turbulent flow to identify the
dominating stress in different regions of the
configuration. \citet{longest2013aerodynamic} performed RANS
predictions of airflow passages with varying flow conditions while
tracking point particles representing agglomerates. The scope was to
evaluate the potential effect of several flow quantities on the
deagglomeration of powders in pharmaceutical inhalers. The evaluated
quantities include the turbulent kinetic energy and the dissipation
rate, the shear stress, and the number of wall hits among others. The
possibility of breakage was evaluated based on a comparison with
experiments conducted in parallel.

Apart from powder inhalers, the particle trajectory modeling approach
has been applied in more general Euler--Lagrange investigations.
\citet{babler2012breakup} conducted a DNS of a diluted particle-laden
homogeneous and isotropic turbulent flow to compute the so-called
exit-time, which describes the time interval needed for agglomerates
to reach flow regions where they become for the first time exposed to
turbulent stresses greater than a critical threshold.  The exit-time
has been used to develop breakup rates for each case, which can be
further used in population balances. In a subsequent study by
\citet{babler2015}, the critical stress was characterized by the size
and fractal dimension of the agglomerate in a power-law relationship,
and the same method was applied to derive breakup rates for various
bounded and unbounded flows. \citet{marchioli2015turbulent} followed a
similar approach to compute the exit-time, arguing that agglomerates
tend to break in a ductile mode due to the accumulation of the
stresses over time rather than the brittle (instantaneous) mechanism
usually assumed. It was postulated that the process of breakage starts
when the acting stress reaches a certain value and ends when the
accumulated amount exceeds another prescribed threshold.

It can be concluded that Euler--Lagrange methods dealing with the
idealized sphere concept have mainly focused on describing the onset
of breakage. Models for the resulting fragment size distribution and
the velocity of the fragments are not well established.  For example,
\citet{Njobuenwu2018} published an investigation on the agglomeration
and breakup of agglomerates in a turbulent channel flow adopting the
breakup model suggested by \citet{babler2015}. The breakup event was
realized by replacing the broken agglomerate by two equally-sized
fragments (symmetrical breakup). However, elaborations on the
post-breakup kinetics of the fragments are missing. \newline

In summary, different stress mechanisms (i.e., turbulent, drag and
rotation) were identified in the literature to take part in the
breakup of agglomerates in fluid flows. In addition, models describing
the strength of agglomerates are also available. That enabled the
specification of breakup conditions in various studies.  However,
refined methodologies which deliver a complete description especially
for the behavior of the arising fragments due to breakage in the
context of the effective sphere approach are still missing.
%
%In summary, the Eulerian--Lagrangian description of particle-laden
%flows has certain advantages in comparison to the Eulerian--Eulerian
%approach
%owing to its ability to track agglomerates through the flow field and
%to investigate the effect of the flow on the breakup process without
%predefined breakup rates. Various methodologies exist within the
%Eulerian--Lagrangian approach to study the breakage phenomena.
%Methods based
%on determining inter-particle contact forces and torques are more
%detailed than those describing agglomerates as porous nearly
%spherical
%structures. For instance, the size and velocity of the fragments
%after
%breakup can be directly predicted. However, due to their high
%computational costs these methods are still restricted to a small
%number of agglomerate comprising a moderate number of primary
%particles. In addition, many simplifying assumptions have to be
%incorporated to ease the handling of the complicated system. The
%brief
%literature overview shows a certain lack of comprehensive numerical
%studies adopting the idealized sphere technique for studying breakup
%processes in turbulent flows, and simultaneously taking breakup and
%agglomeration of a huge number of particles into account.
%

%=====================================================================
%=====================================================================
\section{Wall-Impact Breakage}
\label{sec:litrev_wallimpact}
\fancyhead[OL]{Wall-Impact Breakage}
%=====================================================================
%=====================================================================

In practical applications of particle-laden flow, the interactions
between particles and the surrounding walls are often inevitable,
which renders the mechanism of breakage by wall impact of particular
interest.

In an analogy to the case of the breakup by fluid stresses,
particle-laden flow simulation methodologies which offer forecasting
for the wall-impact breakage phenomena can be distinguished into two
categories based on the level of captured details.  The first
conserves the structure and applies DEM to describe the dynamics of
the primary particles within the agglomerate. If an agglomerate hits a
wall, the contact forces and torques exerted by the wall onto the
wall-touching particles are computed according to the soft-sphere
treatment. Consequently, these particles start to bounce against the
wall leading to a propagation of the impact force into the next layer
of neighboring particles and so forth. As typical for DEM, the
disintegration of particles is directly governed by the force and
torque balances applied on the level of the single
particles. Exemplary studies adopting this approach, which combines
computational fluid dynamics and DEM will be discussed in
Section~\ref{sec:break_wall_DEM}.

In the second method agglomerates are represented by effective
spheres. In the absence of inter-particle contacts, the progress of
the impact energy within the agglomerate causing the breakup of
certain particle contacts cannot be determined. For this purpose,
various efforts have been undertaken in the literature towards
developing models describing the outcome of an impact event based on
the conditions of the impact incident. The contributions relying on
this concept are addressed in Section~\ref{sec:break_wall_replace}.

%=====================================================================
\subsection{Wall-Impact Breakage Predicted by Condition-Free
Techniques}
\label{sec:break_wall_DEM}
%=====================================================================

Investigations employing a direct wall-impact breakage forecasting
have been mostly concerned with dry powder
inhalers. \citet{tong2011numerical} carried out Euler--Lagrange
simulations based on DEM to assess the deagglomeration performance of
different inhaler designs taking mannitol (drug-only) agglomerates
into account. The simulations employed single agglomerates consisting
of 3000 polydisperse particles. Various configurations including five
pipe bend designs and three volumetric flow rates were
investigated. The pipe bends are distinguished by the number of bends
(single or double) and the bend angle.  It was demonstrated that a
pipe bend design that deflects the flow by a 90$^\circ$ using two
45$^\circ$ bends delivers the best performance, i.e., the finest
particle distribution at the outlet and the smallest amount of
particle deposition on the wall. The results of the simulations are
found to be in qualitative agreement with earlier experimental
measurements \citep{adi2010impact}.  \citet{tong2013cfd} applied the
same Euler--Lagrange DEM simulation technique to investigate the
performance of a commercial dry powder inhaler
(Aerolizer{\textregistered}). Eight agglomerates each comprising 3000
mannitol particles were simultaneously taken into account. It was
found that the wall-impact breakage is the most effective
deagglomeration mechanism in the device.

Again, the condition-free prediction of breakage based on DEM is
CPU-time intensive and is currently limited to elementary
investigation and a low number of considered agglomerates.
% In a series on the deagglomeration performance of a commercial dry
% powder inhaler (Aerolizer{\textregistered})
% \citet{coates2004effect} investigated the effect of grid structure
% and
% mouthpiece length, the mouthpiece  inlet size
% \citep{coates2006effect}, and mouthpiece geometry
% \citep{coates2007influence} on the breakage of drug-only
% agglomerates of mannitol
% powders.

%=====================================================================
\subsection{Wall-Impact Breakage Predicted by Condition-Based
Techniques}
\label{sec:break_wall_replace}
%=====================================================================

Efforts towards understanding the wall-impact breakage phenomenon and
developing predictive models have been reported in numerous
experimental works
\citep[e.g.,][]{john1993breakup,froeschke2003impact,samimi2004analysis,davies1951impingement}.
In addition, fundamental DEM studies have been extensively carried out
to study the breakage phenomenon focusing on impact events of single
agglomerates in isolated environments. By precisely controlling the
impact conditions, these simulations allowed to systematically
evaluate the effect of the operating conditions on the fragmentation
providing important insights which are not accessible by experiments
\citep[e.g.,][]{thornton1996numerical,ning1997distinct,mishra2001impact,
  moreno2003effect,tong2009numerical}.  In general, it can be
concluded that the breakage behavior of agglomerates is influenced by
several parameters: The properties of the primary particles and the
impacted wall, the nature of the inter-particle binding mechanism, the
structure of the agglomerate, and the dynamics of the agglomerate,
i.e., the impact velocity and angle \citep{ghadiri2007analysis}. 
\newline

In an early attempt to formulate a criterion for the wall-impact
breakage, \citet{kousaka1979dispersion} suggested a simple model based
on theoretical arguments. They compared the compressive stress at the
center plane of a perfect sphere impacting a wall with the tensile
strength of the agglomerate as suggested by \citet{rumpf1962}.
However, this approach suffers from a number of pitfalls. First, the
stress is linearly related to the so-called collision time, which was
not properly defined. In addition, the formula of the compressive
stress is only valid for normal wall impacts and an extension towards
oblique and shear impact events is not discussed. Lastly, the model
does not offer any information on the resulting fragment size
distribution or the kinetics of the arising fragments.  \newline 

More recent research has followed a different strategy, which uses the
data obtained from experiments and DEM simulations to develop and tune
(semi-)empirical breakage models. A common practice in these studies
is to characterize the outcomes of impact events by means of
dimensionless numbers describing the functional relationship between
various process parameters in a physically meaningful manner.
Subsequently, phenomenological models are proposed in the form of
correlations between the dimensionless numbers and the investigated
breakage measures.

For example, \citet{kafui1993powders} found that the damage ratio
defined as the number of broken contacts over the number of initial
contacts within an agglomerate can be expressed as a function of the
Weber number defined by the impact velocity $v_{\text{imp}}$ and the
diameter $d_\text{pp}$ and the surface energy $\gamma_s$ of the
primary particle, i.e., $\text{We} \, = \, \rho_\text{f} \,
d_{\text{pp}} \, v_{\text{imp}} /(2 \gamma_s)$.  This means that the
damage ratio curves obtained under different conditions tend to
collapse into a single curve when plotted as a function of
We. However, this conclusion was found considering solely normal
impact events carried out with different impact velocities and varying
particle surface energies while keeping all other conditions fixed. In
subsequent studies, \citet{thornton1996numerical} and
\citet{subero1999effect} followed the same approach and suggested that
a better representation of the results can be obtained based on a
modified Weber number, in which the minimum velocity needed for
starting breakage is subtracted from the impact velocity.

\citet{moreno2006mechanistic} derived a new dimensionless number based
on an energy balance assuming that the work needed for breaking all
contacts within an agglomerate is proportional to the incident kinetic
energy. The new dimensionless number can be seen as a product of the
Weber number and a second dimensionless number including the particle
diameter, the modulus of elasticity and the surface energy.  The
proposed dimensionless number allowed an improved unification of the
results, which was achieved over a wider range of surface energies.

\citet{le2011aggregate} carried out 2D simulations for agglomerate
breakage under normal and shear (wall-parallel) impact scenarios.  For
normal impact, a dimensionless number analogous to the Weber number
was considered. However, in the shear impact case a different
dimensionless number was proposed, which describes the ratio between
the frictional work and the work of cohesion. Based on these
considerations, the values of the damage ratio obtained in each case
were well unified by the corresponding dimensionless numbers. 
\newline 

Although these studies are very insightful contributions, it is
important to point out the limitations that prevent a direct
application of the models presented in these studies in particle-laden
flow simulations. First, the investigations strictly focused on
evaluating the effect of cohesion on the breakage of contacts. In
other words, the relationship between a certain dimensionless number
and the damage ratio was investigated while varying the inter-particle
bond strength (i.e., surface energy or Hamaker constant).  However,
the influence of other impact conditions such as the impact angle and
the diameters and number of the primary particles was not
appropriately taken into account since these parameters were fixed in
the simulations. Another shortcoming is that the measure used for
quantifying the results, i.e., the damage ratio, does not allow a
unique description of the post-breakage particle size distribution
since different breakage patterns can be observed for the same damage
ratio \citep{tong2009numerical}. Moreover, the models also lack a
description for the post-breakage motion of the disintegrated
fragments.

In a recent contribution, \citet{van2020analysis} have been probably
the first to present a fragmentation model applicable for
particle-laden flow simulations avoiding some of the main drawbacks
mentioned above. The model relies on the data obtained by DEM
simulations of normal wall-impact events accounting for agglomerates
consisting of monodisperse particles. The reference simulations took
into account a range of impact velocities, particle surface energy
magnitudes, primary particle diameters, and agglomerate sizes.  The
analysis of the results focused on characterizing the fraction of the
detached primary particles based on an extended version of the
dimensionless number proposed by \citet{moreno2006mechanistic}. In
addition, the size distribution of the arising fragments was estimated
by dividing the remaining volume of the original agglomerates into the
largest possible number of discrete size classes after subtracting the
volume of the detached primary particles. The study suggests to
describe the post-breakage velocities of the fragments by means of a
momentum factor. This factor expresses the fraction of the kinetic
energy of the original agglomerate which is devoted to move the
fragments after breakage. Although the introduction of a model that
allows to determine the momentum factor based on the conditions of the
impact event was part of the objectives of the study, a further
elaboration on this issue is not available. Furthermore, it is not
described in which directions the fragments are supposed to spread
after breakage. The main drawback of the model is that it only
provides forecasts for normal impact events since the effect of the
impact angle was not investigated.

It is worth mentioning that the concept of using detailed methods to
develop practical wall-impact breakage models for Euler--Lagrange
point-particle computations has been also applied in the context of
carrier-based agglomerates
\citep[e.g.,][]{vanwachem2017simulation,ariane2018wall}.  However, the
suggested breakage models for carrier-based agglomerates mainly
incorporate arguments specific for particles connected to a larger
carrier.  Thus, they are generally not applicable for traditional
agglomerates. Since the carrier-based formulation of agglomerates is
not of interest for the present thesis as mentioned earlier, such
studies are not further discussed here. The interested reader is
referred to the relevant reviews, e.g.,
\citet{zheng2021flow,yang2015numerical,sommerfeld2019potential}.
\newline

To conclude, models describing the wall-impact breakage in the context
of the description of agglomerates as effective spheres are scarce.
The advantages of the recent data-driven model by
\citet{van2020analysis} over the simple theoretical model by
\citet{kousaka1979dispersion} are clear since the former mimics both
important features of the breakage scenario, i.e., the breakage onset
and the fragment size distribution as predicted in detailed DEM
simulations. However, since solely normal impact events were
considered in the model by \citet{van2020analysis}, the applicability
of the model in flow cases dominated by very small wall-impact angles
must be explored.
%Building on this argument, the model s
% \citet{vanwachem2017simulation} correlated data from number of
% resolved simulations to devise
%the rate of detachment of fine particles from a carrier. The
%detailed
%simulations were designed to study the
%detachment of the fine particles from a single carrier as a results
%of a collision of two carriers, the impact of a carrier at the wall,
%and the acceleration of a carrier in a flow field. In addition, the
%re-attachment of the fines to the carrier was considered.   The
%obtained models were used in LES simulations of complete dry powder
%inhaler imposing a fluid flow obtained from a real inhalation
%profile of a patient. The carrier particles in the inhaler
%simulations were modeled as Lagrangian point masses applying DEM.
%The evolution of the fine particles in the inhaler was taken into
%account by introducing the volumetric concentration of the fines as
%a passive scalar for which a conservation equation is solved. The
%sources and sinks in the scalar equation are described by the
%correlations obtained from the fully-resolved simulations.
%
%\citet{cui2018application} carried out fully-resolved
%(i.e., particle-resolved fluid interactions coupled with DEM)
%lattice-Boltzmann simulations for a carrier particle covered by a
%single layer of monodisperse primary particle exposed to a plug
%flow. The
%detailed results obtained by this approach enabled the derivation of
%simplified analytic models for the detachment of drug particles due
%to
%fluid forces. In a later study,
%\citet{ariane2018wall} performed DEM simulations of
%single carrier-based agglomerate impacting a wall and used the
%results to propose an
%analytic model for the
%detachments of fines from carriers. Note that, in
%general, the detachment models developed for the carrier-based
%agglomerates are based on arguments specific for their own problem.
%Thus they are generally not helpful for predicting the breakage of
%typical agglomerates.
